	\renewcommand{\labelenumi}{\alph{enumi}.)}
	\renewcommand{\labelenumii}{\arabic*.}
{\scriptsize\textbf{\textit{Der 8.\,Kyu wird in aller Regel gemeinsam mit dem 9.\,Kyu absolviert, daher ist für diese Prüfung die doppelte Prüfungsgebühr, für entsprechend 2 Prüfungsmarken, zu entrichten.}}\newline Anmerkung zum \textit{(Chisai-no-)}Mawashi-Geri: in der Chisai-Form wird diese Technik als kleiner Halbkreisfußtritt ohne Drehung des Standbeins ausgeführt und ist technisch vollkommen ausreichend. Die sportliche Variante kann und darf aber auch trainiert und gezeigt werden.\newline
 	\begin{tabularx}{\textwidth}{lX}
		\textit{Vorausgesetzte Techniken}:& Zuki, Age-Uke, Yoko-Uke, Harai-Otoshi-Uke, Soto-Uke, Mae-Geri, \textit{(Chisai-no-)}Mawashi-Geri
	\end{tabularx}\\}
	\begin{enumerate}[nosep]
		\item \textit{\textbf{Kihon-Ido}}:
		\begin{enumerate}[nosep]
			\item Harai-Otoshi-Uke/Gyaku-Zuki Ch\={u}dan\tabto{280pt}(Zenkutsu-Dachi)
			\item Mae-Geri Ch\={u}dan\tabto{280pt}(Zenkutsu-Dachi)
			\item \textit{(Chisai-no-)}Mawashi-Geri Ch\={u}dan\tabto{280pt}(Zenkutsu-Dachi)
			\item Oi-Zuki J\={o}dan/Gyaku-Zuki Ch\={u}dan\tabto{280pt}(Zenkutsu-Dachi)
			\item Soto-Uke Ch\={u}dan/Gyaku-Zuki Ch\={u}dan\tabto{280pt}(Zenkutsu-Dachi)
			\item Age-Uke/Gyaku-Zuki J\={o}dan\tabto{280pt}(Sanchin-Dachi)
			\item Yoko-Uke/Gyaku-Zuki Ch\={u}dan\tabto{280pt}(Sanchin-Dachi)
		\end{enumerate}
		\item \textit{\textbf{Kata}}:
			\begin{enumerate}[nosep]
			\item Taikyoku J\={o}dan
			\item Taikyoku Ch\={u}dan
			\end{enumerate}
		\item \textit{\textbf{Partnerformen}}:
		\begin{enumerate}[nosep]
			\item Kihon-Ippon Kumite:\tabto{280pt}J\={o}dan\&Ch\={u}dan (ggf. Gedan)
		\end{enumerate}
	\end{enumerate}